\section{Conclusion}
\label{sec:conclusion}

We presented \methodname{}, a buffer-free framework for cross-domain multimodal class-incremental learning on HSI+LiDAR data.
Through systematic drift analysis across six fusion architectures, we identified an \emph{architecture-dependent drift asymmetry}---a structural prior of independent-branch designs where spatial features drift significantly more than spectral features.
Exploiting this prior, \methodname{} freezes the spectral branch as a stable anchor and applies asymmetric LoRA adapters to the spatial branches with rank matched to each branch's effective adaptation dimensionality.
Combined with SHINE domain whitening and DCR prototype correction, the framework achieves competitive accuracy ($84.5\%$) with replay methods on a 9-task cross-domain benchmark while requiring ${\sim}75\times$ less extra storage than exemplar buffers.
\methodname{} leverages within-domain training data accessibility---a natural property of fixed-sensor deployments---for adapter fine-tuning and prototype recomputation, eliminating the need for a separate exemplar replay buffer.
A strictly exemplar-free variant degrades catastrophically ($33.9\%$), confirming that within-domain data re-access is architecturally essential for prototype--adapter alignment in multi-block LoRA systems.

\paragraph{Limitations.}
The method's effectiveness relies on backbone designs with spectral--spatial separation---a deliberate choice validated by the observation study, but one that restricts generalization to architectures where spectral and spatial streams are not tightly coupled.
The warm-up length is set heuristically; an automatic convergence criterion would improve robustness.
The 6-ordering evaluation (Table~\ref{tab:domain_order}) confirms consistent advantages over replay baselines (winning 5 of 6 orderings), though with variation of $\pm 2.7$\,pp depending on the initial domain's data volume.
Domain-adaptive warm-up strategies could further reduce this sensitivity.

\paragraph{Reproducibility.}
All experiment result artifacts (per-seed JSONs, per-class CSVs), figure generation scripts, and the complete paper source are included in the supplementary material.
The experiment code, trained checkpoints, and instructions for reproducing results on the three HSI+LiDAR benchmarks will be released upon acceptance.

\paragraph{Future work.}
Promising extensions include: (i)~additional sensor modalities (SAR, multispectral) and longer task sequences; (ii)~domain-adaptive warm-up based on drift convergence criteria; and (iii)~applying the drift asymmetry insight to multi-temporal change detection where temporal modality stability varies.
